\documentclass[10pt,a4paper]{amsart}
\usepackage[margin=19mm]{geometry}
\usepackage{amsmath,amssymb,mathtools}
\usepackage[scr=rsfs,cal=euler]{mathalfa}
\usepackage{xfrac}
\usepackage[unicode,hidelinks,bookmarks=false]{hyperref}
\newcommand{\ie}{i.e.\ }
\newcommand{\cf}{cf.\ }
\newcommand{\resp}{respectively\ }
\newcommand{\Subsection}{\subsection}
\providecommand{\nobreakdash}{\nobreak\hbox{-}}
\setlength{\emergencystretch}{2em}
\multlinegap0pt
\usepackage{bm}
\usepackage{bbm}
\usepackage{dsfont}
\usepackage{xspace}
\usepackage{cancel}
\usepackage{mathtools}
\usepackage{amsthm}
\makeatletter
\@ifundefined{lemma}{\newtheorem{lemma}{Lemma}}{}
\makeatother
\newcommand{\E}{\mathbb{E}}
\newcommand{\R}{\mathbb{R}}
\newcommand{\bA}{\mathbf{A}}
\newcommand{\bB}{\mathbf{B}}
\newcommand{\bD}{\mathbf{D}}
\newcommand{\bGamma}{\bm\Gamma}
\newcommand{\bH}{\mathbf{H}}
\newcommand{\bI}{\mathbf{I}}
\newcommand{\bK}{\mathbf{K}}
\newcommand{\bL}{\mathbf{L}}
\newcommand{\bSigma}{\mathbf{\Sigma}}
\newcommand{\bTheta}{\bm\Theta}
\newcommand{\bV}{\mathbf{V}}
\newcommand{\bX}{\mathbf{X}}
\newcommand{\bbeta}{\bm\beta}
\newcommand{\be}{\mathbf{e}}
\newcommand{\bgamma}{\bm\gamma}
\newcommand{\bmu}{\bm\mu}
\newcommand{\bs}{\mathbf{s}}
\newcommand{\btheta}{\bm\theta}
\newcommand{\bu}{\mathbf{u}}
\newcommand{\bxi}{\bm\xi}
\newcommand{\by}{\mathbf{y}}
\newcommand{\bz}{\mathbf{z}}
\newcommand{\dd}{\mathrm{d}}
\newcommand{\tr}{\mathrm{tr}}
\newcommand{\vc}{\mathrm{vec}}
\title[Multivariate MM-estimators with auxiliary Scale for Linear M]{Multivariate MM-estimators with auxiliary Scale for Linear Models with Structured Covariance Matrices}
\author{Hendrik Paul Lopuhaa}
\date{Source: arXiv:2511.05134v1 (CC BY 4.0)}
\begin{document}
\maketitle

\noindent\emph{Source and licence.} Multivariate MM-estimators with auxiliary Scale for Linear Models with Structured Covariance Matrices. Version arXiv:2511.05134v1 (CC BY 4.0), distributed under the Creative Commons Attribution 4.0 International licence, as recorded in the arXiv licence metadata for this exact version and shown on the abstract page https://arxiv.org/abs/2511.05134v1. Full rights notice: licenses/arxiv-2511.05134-CC-BY-4.0.txt.\par\smallskip

\noindent\textit{Editorial scope.} Counted unit: the Lemma whose source label is \texttt{lem:Lambda derivative} in the article (source0.tex lines 4424--4480, source label \texttt{lem:Lambda derivative}), delivered together with the article's own complete original proof of it (source0.tex lines 4481--5310, one \texttt{proof} environment of 830 lines). No mathematics is written, repaired or re-derived: statement and proof are the article's own text line for line. Required context retained uncounted: def:sigma (lines 488--499); def:Psi (lines 1164--1180); def:V linear (lines 1217--1220); def:L (lines 1228--1238); def:Psi linear (lines 1243--1262); def:PsiV (lines 1264--1277); def:Lambda (lines 1343--1348); def:Lambda-beta-gamma (lines 1351--1364); eq:elliptical (lines 1441--1451); def:alpha1-gamma1 (lines 1472--1494); eq:trace vec (lines 3239--3242); eq:prop K (lines 3977--3980); lem:Lemma 5.1 (lines 3981--3998); lem:change order int diff Lambda (lines 4345--4365). Every definition, display and label of those lines is retained unchanged. Omitted from this package and not redistributed: 1--487, 500--1163, 1181--1216, 1221--1227, 1239--1242, 1263--1263, 1278--1342, 1349--1350, 1365--1440, 1452--1471, 1495--3238, 3243--3976, 3999--4344, 4366--4423, 5311--7645. Nothing inside a retained excerpt is omitted or altered: every retained body is delivered exactly as the article's lines are written, so no omission receipt of any kind accompanies this package. Citations: the retained excerpts carry no citation command at all, so no bibliography entry of the article is transcribed and no reference is redistributed. Reference adaptation: none was needed and none was made; every reference inside the delivered unit resolves inside it, so no unresolved reference or citation is produced. Printed slips kept literal: no printed word, symbol, hypothesis or number of the counted statement or of its proof was altered. Layout and class: the article's own class is \texttt{article} with packages including fontenc, amsthm, amsmath, bm, hyperref, amssymb, a4wide, latexsym, dsfont, color, verbatim, cancel, authblk, appendix; the wrapper is the lane's pinned \texttt{amsart} editorial wrapper at 10pt on A4, which restates the theorem-like environments the retained text needs and adds the article's own macro definitions that the retained text needs (\texttt{\textbackslash E}, \texttt{\textbackslash R}, \texttt{\textbackslash bA}, \texttt{\textbackslash bB}, \texttt{\textbackslash bD}, \texttt{\textbackslash bGamma}, \texttt{\textbackslash bH}, \texttt{\textbackslash bI}, \texttt{\textbackslash bK}, \texttt{\textbackslash bL}, \texttt{\textbackslash bSigma}, \texttt{\textbackslash bTheta}, \texttt{\textbackslash bV}, \texttt{\textbackslash bX}, \texttt{\textbackslash bbeta}, \texttt{\textbackslash be}, \texttt{\textbackslash bgamma}, \texttt{\textbackslash bmu}, \texttt{\textbackslash bs}, \texttt{\textbackslash btheta}, \texttt{\textbackslash bu}, \texttt{\textbackslash bxi}, \texttt{\textbackslash by}, \texttt{\textbackslash bz}, \texttt{\textbackslash dd}, \texttt{\textbackslash tr}, \texttt{\textbackslash vc}), with extra wrapper packages (bm, bbm, dsfont, xspace, cancel, mathtools). The delivered numbering is the unit's own. Assets: no retained span contains a figure environment, a graphics-inclusion command, a PDF-page inclusion, a table or a TikZ picture; no image file is copied into this package, the package declares no asset dependency and no image file is redistributed with it. Licence: version arXiv:2511.05134v1 of this article is distributed under the Creative Commons Attribution 4.0 International licence, as recorded in the arXiv licence metadata for this exact version and shown on the abstract page https://arxiv.org/abs/2511.05134v1. A full rights notice is delivered at licenses/arxiv-2511.05134-CC-BY-4.0.txt. Characters: every retained excerpt is pure ASCII, so the delivered engine, XeLaTeX with its default Latin Modern text font, reports no missing character.
\begin{equation}
\label{def:sigma}
\int
\rho_0
\left(
\frac{\displaystyle\sqrt{(\by-\bX\bbeta_0(P))^T\bGamma(\btheta_0(P))^{-1}(\by-\bX\bbeta_0(P))}}{\sigma}
\right)
\,
\text{d}P(\bs)
=
b_0,
\end{equation}

\begin{equation}
\label{def:Psi}
\begin{split}
\Psi_{\bbeta}(\bs,\bxi,\sigma)
&=
u_1\left(\frac{d}{\sigma}\right)
\bX^T\bV^{-1}(\by-\bX\bbeta)\\
\Psi_{\bgamma,j}(\bs,\bxi,\sigma)
&=
u_1\left(\frac{d}{\sigma}\right)
(\by-\bX\bbeta)^T\bV^{-1}
\bH_{1,j}
\bV^{-1}(\by-\bX\bbeta)\\
&\qquad\qquad\qquad\qquad\quad-
\mathrm{tr}\left(\bV^{-1}\frac{\partial \bV}{\partial \gamma_j}\right)\log|\bV|,
\end{split}
\end{equation}

\begin{equation}
\label{def:V linear}
\bV(\bgamma)=\sum_{j=1}^l\gamma_j\bL_j,
\end{equation}

\begin{equation}
\label{def:L}
\bL
=
\Big[
  \begin{array}{ccc}
    \vc\left(\bL_1\right)
 & \cdots &     \vc\left(\bL_l\right) \\
  \end{array}
\Big].
\end{equation}

\begin{equation}
\label{def:Psi linear}
\begin{split}
\Psi_{\bbeta}(\bs,\bxi,\sigma)
&=
u_1\left(\frac{d}{\sigma}\right)
\bX^T\bV^{-1}(\by-\bX\bbeta)\\
\Psi_{\bgamma}(\bs,\bxi,\sigma)
&=
-\bL^T
\left(
\bV^{-1}
\otimes
\bV^{-1}
\right)
\vc\left(
\Psi_\bV(\bs,\bxi,\sigma)
\right),
\end{split}
\end{equation}

\begin{equation}
\label{def:PsiV}
\Psi_\bV(\bs,\bxi,\sigma)
=
k
u_1\left(\frac{d}{\sigma}\right)
(\by-\bX\bbeta)(\by-\bX\bbeta)^T
-
v_1\left(\frac{d}{\sigma}\right)
\sigma^2\bV
-
\bV
\log|\bV|,
\end{equation}

\begin{equation}
\label{def:Lambda}
\Lambda(\bxi,\sigma)
=
\int \Psi(\bs,\bxi,\sigma)\,\dd P(\bs),
\end{equation}

\begin{equation}
\label{def:Lambda-beta-gamma}
\begin{split}
\Lambda_{\bbeta}(\bxi,\sigma)
&=
\int
\Psi_{\bbeta}(\bs,\bxi,\sigma)\,\dd P(\bs),
\qquad
\Lambda_{\bgamma}(\bxi,\sigma)
=
\int
\Psi_{\bgamma}(\bs,\bxi,\sigma)\,\dd P(\bs).
\end{split}
\end{equation}

\begin{equation}
\label{eq:elliptical}
f_{\bmu,\bSigma}(\by)
=
|\bSigma|^{-1/2}
m\left(
(\by-\bmu)^T
\bSigma^{-1}
(\by-\bmu)
\right),
\end{equation}

\begin{equation}
\label{def:alpha1-gamma1}
\begin{split}
\alpha_1
&=
\mathbb{E}_{\mathbf{0},\bI_k}
\left[
\left(
1-\frac{1}{k}
\right)
\frac{\rho_1'\left(c_\sigma\|\bz\|\right)}{c_\sigma\|\bz\|}
+
\frac1k
\rho_1''\left(c_\sigma\|\bz\|\right)
\right],\\
\gamma_1
&=
\frac{\mathbb{E}_{0,\mathbf{I}_k}
\left[
\rho_1''(c_\sigma\|\bz\|)(c_\sigma\|\bz\|)^2+(k+1)\rho_1'(c_\sigma\|\bz\|)c_\sigma\|\bz\|
\right]}{k+2},
\end{split}
\end{equation}

\begin{equation}
\label{eq:trace vec}
\tr(\bA^T\bB)=\vc(\bA)^T\vc(\bB)
\end{equation}

\begin{equation}
\label{eq:prop K}
\bK_{k,k}\vc(\bA)=\vc(\bA^T).
\end{equation}

\begin{lemma}
\label{lem:Lemma 5.1}
Suppose that $\bz$ has a $k$-variate elliptical contoured density
defined in~\eqref{eq:elliptical}, with parameters $\bmu=\mathbf{0}$ and $\bSigma=\bI_k$.
Then $\bu=\bz/\|\bz\|$ is independent of $\|\bz\|$,
has mean zero and covariance matrix $(1/k)\bI_k$.
Furthermore, $\mathbb{E}_{\mathbf{0},\bI_k}[\bu\bu^T\bu]=\mathbf{0}$
and
\[
\mathbb{E}_{\mathbf{0},\bI_k}
\left[
\vc(\bu\bu^T)\vc(\bu\bu^T)^T
\right]
=
\sigma_1(\bI_{k^2}+\mathbf{K}_{k,k})+\sigma_2\vc(\bI_k)\vc(\bI_k)^T,
\]
where $\sigma_1=\sigma_2=(k(k+2))^{-1}$.
\end{lemma}

\begin{lemma}
\label{lem:change order int diff Lambda}
Let $\Lambda$ be defined by~\eqref{def:Lambda} with~$\Psi$ defined in~\eqref{def:Psi}
and let~$\E\|\bX\|^2<\infty$.
Suppose that $\rho_1$ satisfies (R2) and (R5) and $\bV$ satisfies (V5).
Let $\sigma(P)$ be a solution of~\eqref{def:sigma} and let~$\bxi(P)\in \mathfrak{D}$ be a local minimum of~$R_P(\bbeta,\bV(\bgamma))$.
Let $N\subset \R^k\times \bTheta\times(0,\infty)$ be an open neighborhood of~$(\bxi(P),\sigma(P))$.
Then~$\Lambda$ is continuous differentiable at $(\bxi(P),\sigma(P)$ and
for all $(\bxi,\sigma)\in N$,
\[
\frac{\partial\Lambda(\bxi,\sigma)}{\partial \bxi}
=
\int
\frac{\partial\Psi(\bs,\bxi,\sigma)}{\partial \bxi}\,\dd P(\bs)
\quad\text{and}\quad
\frac{\partial\Lambda(\bxi,\sigma)}{\partial \sigma}
=
\int
\frac{\partial\Psi(\bs,\bxi,\sigma)}{\partial \sigma}\,\dd P(\bs).
\]
\end{lemma}

\begin{lemma}
\label{lem:Lambda derivative}
Suppose that $P$ satisfies~(E) for some $(\bbeta^*,\btheta^*)\in\R^q\times\bTheta$
and that $\E\|\bX\|^2<\infty$.
Suppose that $\rho_1$ satisfies (R2) and (R5) and that $\bV$ satisfies (V5) and has
a linear structure~\eqref{def:V linear}.
Let $\sigma(P)$ be the solution of~\eqref{def:sigma} and let $\bxi(P)=(\bbeta_1(P),\bgamma(P))\in \mathfrak{D}$ be
a local minimum of $R_P(\bbeta,\bV(\bgamma))$.
Suppose that $\bbeta_1(P)=\bbeta^*$ and $\bV(\bgamma(P))=\bSigma/|\bSigma|^{1/k}$.
Let $\Lambda$ be defined by~\eqref{def:Lambda} and~\eqref{def:Lambda-beta-gamma} with $\Psi$ 
from~\eqref{def:Psi linear}.
Then 
\[
\bD_{\bxi}
=
\frac{\partial \Lambda(\bxi(P),\sigma(P))}{\partial \bxi}
=
\left(
  \begin{array}{cc}
%\dfrac{\partial \Lambda_{\bbeta}(\bxi(P),\sigma(P))}{\partial \bbeta} & \mathbf{0} \\
\bD_{\bbeta} & \mathbf{0}\\
    \\
%\mathbf{0} & \dfrac{\partial \Lambda_{\bgamma}(\bxi(P),\sigma(P))}{\partial \bgamma} \\
\mathbf{0} & \bD_{\bgamma} \\
  \end{array}
\right),
\]
where
\begin{equation}
\label{def:derivative Lambda beta}
\bD_{\bbeta}=
\frac{\partial\Lambda_{\bbeta}(\bxi(P),\sigma(P))}{\partial \bbeta}
=
-\alpha_1
|\bSigma|^{1/k}
\E\left[
\bX^T\bSigma^{-1}\bX
\right],
\end{equation}
with $\alpha_1$ defined in~\eqref{def:alpha1-gamma1}, 
and
\begin{equation}
\label{def:derivative Lambda gamma}
\bD_{\bgamma}=
\frac{\partial\Lambda_{\bgamma}(\bxi(P),\sigma(P))}{\partial \bgamma}
=
\omega_1\bL^T\left(\bSigma^{-1}\otimes\bSigma^{-1}\right)\bL
-
\omega_2\bL^T
\vc(\bSigma^{-1})
\vc(\bSigma^{-1})^T
\bL,
\end{equation}
where $\bL=\partial\vc(\bV(\bgamma(P)))/\partial\bgamma^T$ is the $k^2\times l$ matrix given in~\eqref{def:L},
$\omega_1=\sigma^2(P)|\bSigma|^{2/k}\gamma_1$ and $\omega_2=\omega_1/k+|\bSigma|^{2/k}$, 
with $\gamma_1$ defined in~\eqref{def:alpha1-gamma1}.
\end{lemma}

\begin{proof}
For convenience, write $\bxi_P=(\bbeta_{1,P},\bgamma_P)=(\bbeta_1(P),\bgamma(P))$,
$\bV_P=\bV(\bgamma(P))$, and $\sigma_P=\sigma(P)$.
Write $\partial \Lambda/\partial\bxi$ as the block matrix
\begin{equation}
\label{def:Lambda derivative}
\frac{\partial \Lambda(\bxi_P,\sigma_P)}{\partial \bxi}
=
\left(
  \begin{array}{cc}
\dfrac{\partial \Lambda_{\bbeta}(\bxi_P,\sigma_P)}{\partial \bbeta} & \dfrac{\partial \Lambda_{\bbeta}(\bxi_P,\sigma_P)}{\partial \bgamma} \\
    \\
\dfrac{\partial \Lambda_{\bgamma}(\bxi_P,\sigma_P)}{\partial \bbeta} & \dfrac{\partial \Lambda_{\bgamma}(\bxi_P,\sigma_P)}{\partial \bgamma} \\
  \end{array}
\right),
\end{equation}
where $\Lambda_{\bbeta}$ and $\Lambda_{\bgamma}$ are defined in~\eqref{def:Lambda-beta-gamma}
with $\Psi_{\bbeta}$ and $\Psi_{\bgamma}$ from~\eqref{def:Psi linear}.
Because $\rho_1$ and $\bV$ satisfy~(R2), (R5), and~(V5), and $\E\|\bX\|^2<\infty$,
according to Lemma~\ref{lem:change order int diff Lambda} we may
we may interchange integration and differentiation  
in~$\partial\Lambda_{\bbeta}/\partial\bgamma$ and $\partial\Lambda_{\bgamma}/\partial\bbeta$.
It can be seen that these are expectations of an odd function of $\by-\bX\bbeta_{1,P}$,
which means that they are equal to zero, as~$\bbeta_{1,P}=\bbeta^*$ is a point of symmetry of $P$.
Therefore
\begin{equation}
\label{eq:block matrix}
\bD_{\bxi}
=
\frac{\partial \Lambda(\bxi_P,\sigma_P)}{\partial \bxi}
=
\left(
  \begin{array}{cc}
\bD_{\bbeta} & \mathbf{0} \\
    \\
\mathbf{0} & \bD_{\bgamma} \\
  \end{array}
\right).
\end{equation}
It remains to determine $\bD_{\bbeta}=\partial\Lambda_{\bbeta}(\bxi_P,\sigma_P)/\partial\bbeta$
and $\bD_{\bgamma}=\partial\Lambda_{\bgamma}(\bxi_P,\sigma_P)/\partial\bgamma$.
According to Lemma~\ref{lem:change order int diff Lambda}, we have that $\Lambda$ is continuous differentiable at $(\bxi_P,\sigma_P)$ and
that we may interchange
integration and differentiation in $\partial\Lambda_{\bbeta}/\partial\bbeta$, where $\Lambda_{\bbeta}$ is defined in~\eqref{def:Lambda-beta-gamma}
with $\Psi_{\bbeta}$ from~\eqref{def:Psi linear}.
We obtain
\begin{equation}
\label{eq:Dbeta elliptical}
\begin{split}
\bD_{\bbeta}
&=
\int
\frac{\partial}{\partial \bbeta}
\Psi_{\bbeta}(\bs,\bxi_P,\sigma_P)\,\dd P(\bs)\\
&=
-\E
\left[
u_1'\left(\frac{d_P}{\sigma_P}\right)
\frac{\bX^T\bV_P^{-1}\be_P
\be_P^T
\bV_P^{-1}\bX}{\sigma_P d_P}
+
u_1\left(\frac{d_P}{\sigma_P}\right)\bX^T\bV_P^{-1}\bX
\right]\\
&=
-\E
\left[
\E
\left[
u_1'\left(\frac{d_P}{\sigma_P}\right)
\frac{\bX^T\bV_P^{-1}\be_P
\be_P^T
\bV_P^{-1}\bX}{\sigma_P d_P}
+
u_1\left(\frac{d_P}{\sigma_P}\right)\bX^T\bV_P^{-1}\bX
\bigg|
\bX
\right]
\right],
\end{split}
\end{equation}
where $d_P^2=\be_P^T\bV_P^{-1}\be_P$ and $\be_P=\by-\bX\bbeta_{1,P}$.
The inner expectation on the right hand side is the conditional expectation of $\by\mid\bX$,
which can be written as
\[
\bX^T\bV_P^{-1/2}
\E
\left[
u_1'\left(\frac{d_P}{\sigma_P}\right)
\frac{\bV_P^{-1/2}\be_P
\be_P^T
\bV_P^{-1/2}}{\sigma_P d_P}
+
u_1\left(\frac{d_P}{\sigma_P}\right)\bI_k
\,\Bigg|\,
\bX
\right]
\bV_P^{-1/2}\bX.
\]
Because $\bbeta_{1,P}=\bbeta^*$ and $\bV_P=\bSigma/|\bSigma|^{1/k}$, the previous expression is equal to
\[
|\bSigma|^{1/k}
\bX^T\bSigma^{-1/2}
\E
\left[
u_1'\left(\frac{d^*}{\sigma_P}\right)
\frac{|\bSigma|^{1/k}\bSigma^{-1/2}\be^*(\be^*)^T\bSigma^{-1/2}}{\sigma_Pd^*}
+
u_1\left(\frac{d^*}{\sigma_P}\right)\bI_k
\,\Bigg|\,
\bX
\right]
\bSigma^{-1/2}\bX,
\]
where $(d^*)^2=|\bSigma|^{1/k}(\by-\bX\bbeta^*)^T\bSigma^{-1}(\by-\bX\bbeta^*)$ and $\be^*=\by-\bX\bbeta^*$.
Note that $\by\mid\bX$ has the same distribution as $\bSigma^{1/2}\bz+\bX\bbeta^*$,
where~$\bz$ has a spherical density $f_{\mathbf{0},\bI_k}$,
so that the expression in the previous display is equal to
\[
|\bSigma|^{1/k}
\bX^T
\bSigma^{-1/2}
\mathbb{E}_{\mathbf{0},\bI_k}
\left[
u_1'\left(c_\sigma\|\bz\|\right)
\frac{c_\sigma}{\|\bz\|}
\bz\bz^T
+
u_1\left(c_\sigma\|\bz\|\right)
\bI_k
\right]
\bSigma^{-1/2}
\bX,
\]
where $c_\sigma=|\bSigma|^{1/(2k)}/\sigma_P$.
Let $\bu=\bz/\|\bz\|$.
Then with Lemma~\ref{lem:Lemma 5.1} we find
\[
\begin{split}
&
\mathbb{E}_{\mathbf{0},\bI_k}
\left[
u_1'\left(c_\sigma\|\bz\|\right)
\frac{c_\sigma}{\|\bz\|}
\bz\bz^T
+
u_1\left(c_\sigma\|\bz\|\right)
\bI_k
\right]\\
&=
\mathbb{E}_{\mathbf{0},\bI_k}
\left[
u_1'\left(c_\sigma\|\bz\|\right)c_\sigma\|\bz\|
\right]
\mathbb{E}_{\mathbf{0},\bI_k}
\left[
\bu\bu^T
\right]
+
\mathbb{E}_{\mathbf{0},\bI_k}
\left[
u_1\left(c_\sigma\|\bz\|\right)
\right]
\bI_k
=
\alpha_1\bI_k,
\end{split}
\]
where
\[
\begin{split}
\alpha_1
&=
\mathbb{E}_{\mathbf{0},\bI_k}
\left[
\frac{1}{k}
u_1'\left(c_\sigma\|\bz\|\right)c_\sigma\|\bz\|
+
u_1\left(c_\sigma\|\bz\|\right)
\right]\\
&=
\mathbb{E}_{\mathbf{0},\bI_k}
\left[
\left(
1-\frac{1}{k}
\right)
\frac{\rho_1'\left(c_\sigma\|\bz\|\right)}{c_\sigma\|\bz\|}
+
\frac1k
\rho_1''\left(c_\sigma\|\bz\|\right)
\right].
\end{split}
\]
We conclude that
\[
\dfrac{\partial \Lambda_{\bbeta}(\bxi_P,\sigma_P)}{\partial \bbeta}
=
-
\alpha_1
|\bSigma|^{1/k}
\E\left[
\bX^T\bSigma^{-1}\bX
\right].
\]
Next, we determine $\partial\Lambda_{\bgamma}(\bxi_P,\sigma_P)/\partial\bgamma$.
From~\eqref{def:Psi linear} we have
\[
\Psi_{\bgamma,j}
=
-
\vc(\bV^{-1}\bL_j\bV^{-1})^T
\vc\left(
\Psi_\bV
\right),
\]
for all $j=1,2,\ldots,l$,
where $\Psi_\bV$ is defined in~\eqref{def:PsiV}. 
Because $|\bV_P|=1$, we have
\[
\begin{split}
\int \Psi_{\bV}(\bs,\bxi_P,\sigma_P)\,\text{d}P(\bs)
&=
\E\left[
k
u_1\left(\frac{d_P}{\sigma_P}\right)
\be_P\be_P^T
-
v_1\left(\frac{d_P}{\sigma_P}\right)
\sigma_P^2\bV_P
\right]\\
&=
\E\left[
\E\left[
k
u_1\left(\frac{d_P}{\sigma_P}\right)
\be_P\be_P^T
-
v_1\left(\frac{d_P}{\sigma_P}\right)
\sigma_P^2\bV_P
\bigg|
\bX
\right]
\right].
\end{split}
\]
Similar to the reasoning before, the inner expectation can be written as
\[
\begin{split}
&
\E_{\mathbf{0},\bI_k}
\left[
ku_1(c_\sigma\|\bz\|)
\bSigma^{1/2}
\bz\bz^T
\bSigma^{1/2}
-
v_1(c_\sigma\|\bz\|)
\sigma_P^2
\bSigma/|\bSigma|^{1/k}
\right]\\
&=
k
\bSigma^{1/2}
\E_{\mathbf{0},\bI_k}
\left[
u_1(c_\sigma\|\bz\|)
\bz\bz^T
\right]
\bSigma^{1/2}
-
\E_{\mathbf{0},\bI_k}
\left[
u_1(c_\sigma\|\bz\|)\|\bz\|^2
\right]
\bSigma
\\
&=
k
\bSigma^{1/2}
\E_{\mathbf{0},\bI_k}
\left[
u_1(c_\sigma\|\bz\|)\|\bz\|^2
\right]
\E_{\mathbf{0},\bI_k}
\left[
\bu\bu^T
\right]
\bSigma^{1/2}
-
\E_{\mathbf{0},\bI_k}
\left[
u_1(c_\sigma\|\bz\|)\|\bz\|^2
\right]
\bSigma\\
&=
\mathbf{0},
\end{split}
\]
since $\E_{\mathbf{0},\bI_k}[\bu\bu^T]=(1/k)\bI_k$, according to Lemma~\ref{lem:Lemma 5.1}.
Hence we conclude that 
\[
\int \Psi_{\bV}(\bs,\bxi_P,\sigma_P)\,\text{d}P(\bs)=\mathbf{0}.
\]
Since, we may interchange integration and differentiation in $\partial\Lambda_{\bgamma}/\partial\bgamma$,
according to Lemma~\ref{lem:change order int diff Lambda},
this means that
for each $j,s=1,\ldots,l$,
\begin{equation}
\label{eq:Dlambda-gamma}
\begin{split}
\frac{\partial\Lambda_{\bgamma,j}(\bxi_P,\sigma_P)}{\partial \gamma_s}
&=
-
\vc\left(\bV_P^{-1}\bL_j\bV_P^{-1}\right)^T\vc\left(
\int
\frac{\partial\Psi_{\bV}(\bs,\bxi_P,\sigma_P)}{\partial \gamma_s}
\,\dd P(\bs)
\right),
\end{split}
\end{equation}
where $\Psi_\bV$ is defined in~\eqref{def:PsiV}.
We have
\begin{equation}
\label{eq:decomp PsiV}
\begin{split}
\frac{\partial \Psi_\bV}{\partial\gamma_s}
&=
\frac{\partial}{\partial\gamma_s}
k
u_1\left(\frac{d}{\sigma}\right)
(\by-\bX\bbeta)(\by-\bX\bbeta)^T
-
\frac{\partial}{\partial\gamma_s}
v_1\left(\frac{d}{\sigma}\right)\sigma^2\mathbf{V}
+
\frac{\partial}{\partial\gamma_s}
(\log|\bV|)\mathbf{V}.
\end{split}
\end{equation}
Because $\bV$ satisfies~\eqref{def:V linear}, if follows that $\partial\bV/\partial\gamma_s=\bL_s$.
Similar to~\eqref{eq:Dbeta elliptical}, for the first term in~\eqref{eq:decomp PsiV} we have at~$(\bxi_P,\sigma_P)$:
\[
\begin{split}
&\int
\frac{\partial}{\partial\gamma_s}
k
u_1\left(\frac{d_P}{\sigma_P}\right)
\be_P\be_P^T
\,\text{d}P(\bs)\\
&\quad=
\E
\left[
\E
\left[
-k
u_1'\left(\frac{d_P}{\sigma_P}\right)
\frac{1}{2\sigma_P d_P}
\be_P^T\bV_P^{-1}\frac{\partial\bV_P}{\partial\gamma_s}\bV_P^{-1}\be_P
\cdot
\be_P\be_P^T
\bigg|
\bX
\right]
\right],
\end{split}
\]
where $d_P^2=\be_P^T\bV_P^{-1}\be_P$ and $\be_P=\by-\bX\bbeta_{1,P}$.
Because $\bbeta_{1,P}=\bbeta^*$ and $\bV_P=\bSigma/|\bSigma|^{1/k}$, 
the inner expectation on the right hand side can be written as
\[
\begin{split}
&
-k
\E
\left[
u_1'\left(\frac{d^*}{\sigma_P}\right)
\frac{|\bSigma|^{2/k}}{2\sigma_Pd^*}
(\be^*)^T\bSigma^{-1}\bL_s\bSigma^{-1}\be^*
\cdot
\be^*(\be^*)^T
\right]\\
&=
-
\sigma_P^2
\E_{\mathbf{0},\bI_k}
\left[
\frac{ku_1'\left(c_\sigma\|\bz\|\right)(c_\sigma\|\bz\|)^3}2
\bu^T\bSigma^{-1/2}\bL_s\bSigma^{-1/2}\bu
\cdot
\bSigma^{1/2}\bu\bu^T\bSigma^{1/2}
\right],
\end{split}
\]
where $\bz=\bSigma^{-1/2}(\by-\bX\bbeta^*)$ and $\bu=\bz/\|\bz\|$.
According to Lemma~\ref{lem:Lemma 5.1}, 
the expectation of the right hand side is equal to
\begin{equation}
\label{eq:term1 decomp PsiV}
-
\sigma_P^2
\E_{\mathbf{0},\bI_k}
\left[
\frac{ku_1'\left(c_\sigma\|\bz\|\right)(c_\sigma\|\bz\|)^3}{2}
\right]
\E_{\mathbf{0},\bI_k}
\left[
\bu^T\bSigma^{-1/2}\bL_s\bSigma^{-1/2}\bu
\bSigma^{1/2}\bu\bu^T\bSigma^{1/2}
\right].
\end{equation}
For the second term in~\eqref{eq:decomp PsiV} we get at $(\bxi_P,\sigma_P)$:
\[
\begin{split}
&
-
\int\frac{\partial}{\partial\gamma_s}
v_1\left(\frac{d_P}{\sigma_P}\right)\sigma_P^2\mathbf{V}_P
\,\text{d}P(\bs)\\
&=
\E\left[
\E\left[
v_1'\left(\frac{d_P}{\sigma_P}\right)
\frac{1}{2\sigma_P d_P}
\be_P^T\bV_P^{-1}\bL_s\bV_P^{-1}\be_P
\cdot\sigma_P^2\mathbf{V}_P
-
v_1\left(\frac{d_P}{\sigma_P}\right)\sigma_P^2\bL_s
\bigg|\bX
\right]
\right],
\end{split}
\]
where $d_P^2=\be_P^T\bV_P^{-1}\be_P$ and $\be_P=\by-\bX\bbeta_{1,P}$.
As before, the inner expectation can be written as
\[
\begin{split}
&
\sigma_P^2
\E\left[
v_1'\left(\frac{d^*}{\sigma_P}\right)
\frac{|\bSigma|^{1/k}}{2\sigma_Pd^*}
(\be^*)^T\bSigma^{-1}\bL_s\bSigma^{-1}\be^*
\cdot
\bSigma
-
v_1\left(\frac{d^*}{\sigma_P}\right)\sigma_P^2
\bL_s
\right]\\
&=
\sigma_P^2
\E_{\mathbf{0},\bI_k}
\left[
\frac{v_1'\left(c_\sigma\|\bz\|\right)c_\sigma\|\bz\|}{2}
\bu^T\bSigma^{-1/2}\bL_s\bSigma^{-1/2}\bu
\cdot
\bSigma
-
\sigma_P^2
v_1\left(c_\sigma\|\bz\|\right)\bL_s
\right].
\end{split}
\]
With Lemma~\ref{lem:Lemma 5.1}, 
for the expectation of the second term in~\eqref{eq:decomp PsiV} we get 
\begin{equation}
\label{eq:term2 decomp PsiV}
\begin{split}
\sigma_P^2
\E_{\mathbf{0},\bI_k}
\left[
\frac{v_1'\left(c_\sigma\|\bz\|\right)c_\sigma\|\bz\|}{2}
\right]
\E_{\mathbf{0},\bI_k}
\left[
\bu^T\bSigma^{-1/2}\bL_s\bSigma^{-1/2}\bu
\right]
\bSigma\\
-
\sigma_P^2
\E_{\mathbf{0},\bI_k}
\left[
v_1\left(c_\sigma\|\bz\|\right)
\right]
\bL_s.
\end{split}
\end{equation}
For the third term in~\eqref{eq:decomp PsiV} we get at $(\bxi_P,\sigma_P)$:
\begin{equation}
\label{eq:term3 decomp PsiV}
\begin{split}
\frac{\partial}{\partial\gamma_s}
(\log|\bV_P|)\mathbf{V}_P
&=
\left(
\frac{\partial\log|\bV_P|}{\partial\gamma_s}
\right)
\bV_P
+
\left(
\log|\bV_P|
\right)
\frac{\partial\bV}{\partial\gamma_s}\\
&=
\text{tr}
\left(
\bV_P^{-1}\bL_s
\right)
\bV_P
=
\text{tr}
\left(
\bSigma^{-1}\bL_s
\right)
\bSigma,
\end{split}
\end{equation}
using that $|\bV_P|=1$.
It follows that 
\begin{equation}
\label{eq:Dlambda-gamma-js}
\begin{split}
&
\int
\frac{\partial\Psi_{\bV}(\bs,\bxi_P,\sigma_P)}{\partial \gamma_s}
\,\dd P(\bs)\\
&=
-
\sigma_P^2
\E_{\mathbf{0},\bI_k}
\left[
\frac{ku_1'\left(c_\sigma\|\bz\|\right)(c_\sigma\|\bz\|)^3}{2}
\right]
\E_{\mathbf{0},\bI_k}
\left[
\bu^T\bSigma^{-1/2}\bL_s\bSigma^{-1/2}\bu
\bSigma^{1/2}\bu\bu^T\bSigma^{1/2}
\right]\\
&\qquad+
\sigma_P^2
\E_{\mathbf{0},\bI_k}
\left[
\frac{v_1'\left(c_\sigma\|\bz\|\right)c_\sigma\|\bz\|}{2}
\right]
\E_{\mathbf{0},\bI_k}
\left[
\bu^T\bSigma^{-1/2}\bL_s\bSigma^{-1/2}\bu
\right]
\bSigma\\
&\qquad\qquad-
\sigma_P^2
\E_{\mathbf{0},\bI_k}
\left[
v_1\left(c_\sigma\|\bz\|\right)
\right]
\bL_s
+
\text{tr}
\left(
\bSigma^{-1}\bL_s
\right)
\bSigma.
\end{split}
\end{equation}
In view of~\eqref{eq:Dlambda-gamma} and~\eqref{eq:Dlambda-gamma-js}, for the first term in $\partial\Lambda_{\bgamma,j}/\partial\gamma_s$
we obtain
\[
\begin{split}
&
\vc(\bV_P^{-1}\bL_j\bV_P^{-1})^T
\vc\left(
\E_{\mathbf{0},\bI_k}
\left[
\bu^T\bSigma^{-1/2}\bL_s\bSigma^{-1/2}\bu
\bSigma^{1/2}\bu\bu^T\bSigma^{1/2}
\right]
\right)\\
&=
|\bSigma|^{2/k}
\vc(\bSigma^{-1}\bL_j\bSigma^{-1})^T
\E_{\mathbf{0},\bI_k}
\left[
\vc\left(
\bSigma^{1/2}\bu\bu^T\bSigma^{1/2}
\right)
\bu^T\bSigma^{-1/2}\bL_s\bSigma^{-1/2}\bu
\right]\\
&=
|\bSigma|^{2/k}
\vc(\bSigma^{-1}\bL_j\bSigma^{-1})^T
(\bSigma^{1/2}\otimes\bSigma^{1/2})
\E_{\mathbf{0},\bI_k}
\left[
\vc\left(
\bu\bu^T
\right)
\vc(\bu\bu^T)^T
\right]\\
&\qquad\qquad\qquad\qquad\qquad\qquad\qquad\qquad\qquad\qquad\qquad
\vc\left(\bSigma^{-1/2}\bL_s\bSigma^{-1/2}\right)
\\
&=
|\bSigma|^{2/k}
\vc(\bSigma^{-1/2}\bL_j\bSigma^{-1/2})^T
\frac{1}{k(k+2)}
\Big(
\bI_{k^2}+\bK_{k,k}
+
\vc(\bI_k)\vc(\bI_k)^T
\Big)\\
&\qquad\qquad\qquad\qquad\qquad\qquad\qquad\qquad\qquad\qquad\qquad
\vc\left(\bSigma^{-1/2}\bL_s\bSigma^{-1/2}\right),
\end{split}
\]
using Lemma~\ref{lem:Lemma 5.1}.
Application of properties~\eqref{eq:prop K} and~\eqref{eq:trace vec}, yields
\[
\begin{split}
&
\vc(\bV_P^{-1}\bL_j\bV_P^{-1})^T
\vc\left(
\E_{\mathbf{0},\bI_k}
\left[
\bu^T\bSigma^{-1/2}\bL_s\bSigma^{-1/2}\bu
\bSigma^{1/2}\bu\bu^T\bSigma^{1/2}
\right]
\right)\\
&=
\frac{|\bSigma|^{2/k}}{k(k+2)}
\Big(
2\text{tr}(\bSigma^{-1}\bL_j\bSigma^{-1}\bL_s)
+
\text{tr}(\bSigma^{-1}\bL_j)\text{tr}(\bSigma^{-1}\bL_s)
\Big).
\end{split}
\]
It follows that the first term in $\partial\Lambda_{\bgamma,j}/\partial\gamma_s$ is equal to
\begin{equation}
\label{eq:first term}
\sigma_P^2
|\bSigma|^{2/k}
\frac{\E_{\mathbf{0},\bI_k}
\left[
u_1'(c_\sigma\|\bz\|)(c_\sigma\|\bz\|)^3
\right]}{2(k+2)}
\Big(
2\text{tr}(\bSigma^{-1}\bL_j\bSigma^{-1}\bL_s)+\text{tr}(\bSigma^{-1}\bL_j)\text{tr}(\bSigma^{-1}\bL_s)
\Big).
\end{equation}
Similarly, for the second term in $\partial\Lambda_{\bgamma,j}/\partial\gamma_s$
we obtain
\[
\begin{split}
&
\vc\left(\bV_P^{-1}\bL_j\bV_P^{-1}\right)^T
\vc\left(\E_{\mathbf{0},\bI_k}
\left[
\bu^T\bSigma^{-1/2}\bL_s\bSigma^{-1/2}\bu
\right]
\bSigma
\right)\\
&=
|\bSigma|^{2/k}
\vc\left(\bSigma^{-1}\bL_j\bSigma^{-1}\right)^T
\text{tr}
\left(
\E_{\mathbf{0},\bI_k}
\left[
\bu\bu^T
\right]
\bSigma^{-1/2}\bL_s\bSigma^{-1/2}
\right)
\vc\left(\bSigma\right)
\\
&=
|\bSigma|^{2/k}
\text{tr}(\bSigma^{-1}\bL_s)
\vc\left(\bSigma^{-1}\bL_j\bSigma^{-1}\right)^T
\vc\left(\bSigma\right)\\
&=
|\bSigma|^{2/k}
\text{tr}(\bSigma^{-1}\bL_s)
\text{tr}(\bSigma^{-1}\bL_j).
\end{split}
\]
It follows that the second term in $\partial\Lambda_{\bgamma,j}/\partial\gamma_s$ is equal to
\begin{equation}
\label{eq:second term}
-\sigma_P^2
|\bSigma|^{2/k}
\frac{\E_{\mathbf{0},\bI_k}
\left[
v_1'(c_\sigma\|\bz\|)c_\sigma\|\bz\|
\right]}{2k}
\text{tr}\left(\bSigma^{-1}\bL_s\right)
\text{tr}\left(\bSigma^{-1}\bL_j\right).
\end{equation}
Finally, the third term in $\partial\Lambda_{\bgamma,j}/\partial\gamma_s$ is equal to
\begin{equation}
\label{eq:third term}
\begin{split}
&
\sigma_P^2
|\bSigma|^{2/k}
\E_{\mathbf{0},\bI_k}
\left[
v_1(c_\sigma\|\bz\|)
\right]
\vc(\bSigma^{-1}\bL_j\bSigma^{-1})^T
\vc(\bL_s)\\
&=
\sigma_P^2
|\bSigma|^{2/k}
\E_{\mathbf{0},\bI_k}
\left[
v_1(c_\sigma\|\bz\|)
\right]
\text{tr}(\bSigma^{-1}\bL_j\bSigma^{-1}\bL_s),
\end{split}
\end{equation}
and the fourth term in $\partial\Lambda_{\bgamma,j}/\partial\gamma_s$ is equal to
\begin{equation}
\label{eq:fourth term}
-|\bSigma|^{2/k}
\text{tr}
\left(
\bSigma^{-1}\bL_s
\right)
\vc(\bSigma^{-1}\bL_j\bSigma^{-1})^T(\vc\bSigma)
=
-|\bSigma|^{2/k}
\text{tr}
\left(
\bSigma^{-1}\bL_s
\right)
\text{tr}
\left(
\bSigma^{-1}\bL_j
\right).
\end{equation}
We conclude that $\partial\Lambda_{\bgamma,j}(\bxi_P)/\partial \gamma_s$ consists of
a term
$\text{tr}(\bSigma^{-1}\bL_j\bSigma^{-1}\bL_s)$
from~\eqref{eq:first term} and~\eqref{eq:third term}
with coefficient
\[
\begin{split}
\omega_1
&=
\sigma_P^2
|\bSigma|^{2/k}
\left(
\frac{\E_{\mathbf{0},\bI_k}
\left[
u_1'(c_\sigma\|\bz\|)(c_\sigma\|\bz\|)^3
\right]}{k+2}
+
\E_{\mathbf{0},\bI_k}
\left[
v_1(c_\sigma\|\bz\|)
\right]
\right)\\
&=
\sigma_P^2
|\bSigma|^{2/k}
\frac{\mathbb{E}_{0,\mathbf{I}_k}
\left[
\rho''(c_\sigma\|\bz\|)(c_\sigma\|\bz\|)^2+(k+1)\rho'(c_\sigma\|\bz\|)c_\sigma\|\bz\|
\right]}{k+2},
\end{split}
\]
and a term $-\text{tr}(\bSigma^{-1}\bL_s)\text{tr}(\bSigma^{-1}\bL_j)$ 
from~\eqref{eq:first term}, \eqref{eq:second term}, and~\eqref{eq:fourth term}
with coefficient
\[
\begin{split}
\omega_2
&=
\sigma_P^2
|\bSigma|^{2/k}
\left(
\frac{\E_{\mathbf{0},\bI_k}
\left[
u_1'(c_\sigma\|\bz\|)(c_\sigma\|\bz\|)^3
\right]}{2(k+2)}
-
\frac{\E_{\mathbf{0},\bI_k}
\left[
v_1'(c_\sigma\|\bz\|)c_\sigma\|\bz\|
\right]}{2k}
\right)
+
|\bSigma|^{2/k}\\
&=
\sigma_P^2
|\bSigma|^{2/k}
\frac{\mathbb{E}_{0,\mathbf{I}_k}
\left[
\rho''(c_\sigma\|\bz\|)(c_\sigma\|\bz\|)^2+(k+1)\rho'(c_\sigma\|\bz\|)c_\sigma\|\bz\|
\right]}{k(k+2)}
+
|\bSigma|^{2/k}\\
&=
\frac{\omega_1}{k}+|\bSigma|^{2/k}.
\end{split}
\]
From the definition of $\bL$ in~\eqref{def:L} it follows that the $l\times l$ matrix with entries
\[
\begin{split}
\frac{\partial\Lambda_{\bgamma,j}(\bxi_P,\sigma_P)}{\partial \gamma_s}
&=
\omega_1\text{tr}(\bSigma^{-1}\bL_j\bSigma^{-1}\bL_s)
-
\omega_2\text{tr}(\bSigma^{-1}\bL_s)\text{tr}(\bSigma^{-1}\bL_j),
\end{split}
\]
is the matrix
\[
\bD_{\bgamma}
=
\frac{\partial\Lambda_{\bgamma}(\bxi_P,\sigma_P)}{\partial \bgamma}
=
\omega_1\mathbf{L}^T\left(\bSigma^{-1}\otimes\bSigma^{-1}\right)\mathbf{L}
-
\omega_2\mathbf{L}^T
\vc(\bSigma^{-1})
\vc(\bSigma^{-1})^T
\mathbf{L}.
\]
This proves the lemma.
\end{proof}

\end{document}
